(1) 左区自由波数 \(k=\sqrt{2mE}/\hbar\)，右区传播波数 \(q=\sqrt{2m(E-V_0)}/\hbar\)。令左方波函数 \(\psi_L=e^{ikx}+r e^{-ikx}\)，右方无向左来波为 \(\psi_R=t e^{iqx}\)。有限势阶处要求 \(\psi_L(0)=\psi_R(0)\) 和 \(\psi_L'(0)=\psi_R'(0)\)，即 \(1+r=t\)、\(k(1-r)=qt\)。
(2) 解联立边界条件得 \(r=(k-q)/(k+q)\)、\(t=2k/(k+q)\)。概率流密度是 \(j=(\hbar/m)\operatorname{Im}(\psi^*\psi')\)，因此左区入射流 \(j_{\rm in}=\hbar k/m\)、反射流大小为 \(|r|^2j_{\rm in}\)，右区透射流 \(j_t=(\hbar q/m)|t|^2\)。反射率 \(R=(k-q)^2/(k+q)^2\)、透射率 \(T=(q/k)|t|^2=4kq/(k+q)^2\)，显然 \(R+T=1\)。
(3) 当 \(E<V_0\) 时，定义 \(\kappa=\sqrt{2m(V_0-E)}/\hbar\)，右区不可归一化的增长指数必须舍去，留下 \(\psi_R=t e^{-\kappa x}\)，穿透长度 \(\ell_p=1/\kappa\)。把 \(q\) 替为 \(i\kappa\) 得 \(r=(k-i\kappa)/(k+i\kappa)\)，其模为1，所以 \(R=1\)。纯实指数的空间概率流 \(j_R=0\)，倏逝概率密度存在却没有穿过半无限势区的持续透射流。
(4) 当 \(E\gg V_0\)，\(q/k=\sqrt{1-V_0/E}=1-V_0/(2E)+O(E^{-2})\)，所以 \(r\simeq V_0/(4E)\)、\(R\simeq V_0^2/(16E^2)\)。当 \(E\downarrow V_0\) 且仍高于势阶时 \(q\to0^+\)，于是 \(T=4kq/(k+q)^2\to0\)、\(R\to1\)；这与低于势阶时全反射的连续极限一致，但不代表右区波函数振幅必为零。
